% Lecture notes: Week 1, Lecture 3 -- Equivalent Load
% To hide this lecture from the website, add a line containing only:  % draft
% To set the date shown on the website, add a line like:  % posted: 2026-09-02
\documentclass[11pt]{article}
\usepackage{mech2030notes}
\begin{document}

% ##################################################################
%  WEEK 1, LECTURE 3
% ##################################################################
\lecture{1}{3}{Equivalent Load}

\section*{Example: Wing Loading on an Airplane}

\noindent
\begin{minipage}[c]{0.56\textwidth}
\centering
\begin{tikzpicture}[scale=0.85]
  % wings (top-down view)
  \draw[thick, fill=gray!8] (0.5,0.9) -- (5.2,-0.15) -- (5.2,-0.45) -- (0.5,-0.9) -- cycle;
  \draw[thick, fill=gray!8] (-0.5,0.9) -- (-5.2,-0.15) -- (-5.2,-0.45) -- (-0.5,-0.9) -- cycle;
  % fuselage
  \draw[thick, fill=gray!8] (0,2.6) to[out=0,in=90] (0.6,1.6) -- (0.6,-2.3)
        to[out=-90,in=0] (0,-2.8) to[out=180,in=-90] (-0.6,-2.3) -- (-0.6,1.6)
        to[out=90,in=180] (0,2.6);
  \com{(0,0)}
  \draw[thinarrow] (-2.4,-1.7) -- (-0.2,-0.12);
  \node[below, font=\small] at (-2.4,-1.7) {center of mass};
  % chord
  \draw[{Stealth[length=1.8mm]}-{Stealth[length=1.8mm]}, thin] (3,0.38) -- (3,-0.68);
  \node[above, font=\small] at (3,0.45) {chord $c(x)$};
  \node[font=\small] at (0,-3.2) {(top-down view)};
\end{tikzpicture}
\end{minipage}%
\hfill
\begin{minipage}[c]{0.4\textwidth}
\begin{itemize}[label=$\rightarrow$]
  \item Distributed load due to lift:
        \[ p_{\text{lift}}(x) \sim c(x) \]
        (over-simplified, but roughly correct)
  \item What is the equivalent load on the plane's center of mass (C.O.M.)?
  \item Will the plane stay in the air?
\end{itemize}
\end{minipage}

\bigskip
\noindent
Lift on one wing:

\begin{center}
\begin{tikzpicture}[baseline=(b), scale=0.8]
  \coordinate (b) at (0,-0.5);
  \draw[thick] (0.35,0.65) to[out=-60,in=60] (0.35,-0.65);
  \draw[thick] (0.65,0) -- (6.5,0);
  \com{(0,0)}
  \draw[thick] (0.65,-1.6) -- (6.5,0);
  \foreach \i in {0,...,7} {
    \pgfmathsetmacro{\xx}{0.65+5.85*\i/8}
    \pgfmathsetmacro{\yy}{-1.6*(1-\i/8)}
    \draw[-{Stealth[length=2mm,width=1.6mm]}, semithick] (\xx,\yy) -- (\xx,-0.03);
  }
  \node[below] at (0.65,-1.6) {$\SI{7000}{lb/ft}$};
  \node[right] at (6.5,-0.25) {$\SI{0}{lb/ft}$};
  \draw[dim] (0,0.9) -- (0.65,0.9) node[midway, above] {\small\SI{10}{ft}};
  \draw[dim] (0.65,0.9) -- (6.5,0.9) node[midway, fill=white] {$\SI{90}{ft}$};
\end{tikzpicture}
\quad{\Large $\equiv$}\quad
\begin{tikzpicture}[baseline=(b), scale=0.8]
  \coordinate (b) at (0,-0.5);
  \draw[thick] (0.35,0.65) to[out=-60,in=60] (0.35,-0.65);
  \draw[thick] (0.65,0) -- (6.5,0);
  \com{(0,0)}
  \draw[force] (2.6,-1.4) node[below] {$F_y^T$} -- (2.6,-0.03);
  \draw[dimto] (0,0.9) -- (2.6,0.9) node[right] {$x_c^T$};
\end{tikzpicture}
\end{center}

\[
  \left\{
  \begin{aligned}
    F_y^T &= \tfrac12\,(\SI{7000}{lb/ft})(\SI{90}{ft}) = \SI{315000}{lb} \\
    x_c^T &= \SI{10}{ft} + \tfrac13(\SI{90}{ft}) = \SI{40}{ft}
  \end{aligned}
  \right.
\]

\subsection*{FBD}

\noindent
\begin{minipage}[b]{0.46\textwidth}
\centering
\begin{tikzpicture}
  \draw[thick] (0,0) -- (4,0);
  \com{(0,0)}
  \draw[force] (0,1.3) node[above] {$W$} -- (0,0.2);
  \draw[force] (2.6,-1.3) node[below] {$F_y^T$} -- (2.6,-0.03);
  \draw[dimto] (0,-0.6) -- (2.55,-0.6);
  \node[below] at (1.3,-0.6) {$x_c^T$};
\end{tikzpicture}\\[0.4em]
single force at $x_c^T$
\end{minipage}%
\hfill{\Large $\equiv$}\hfill
\begin{minipage}[b]{0.46\textwidth}
\centering
\begin{tikzpicture}
  \draw[thick] (0,0) -- (4,0);
  \com{(0,0)}
  \draw[force] (0,1.3) node[above] {$W$} -- (0,0.2);
  \draw[force] (0,-1.3) node[below] {$F_y^T$} -- (0,-0.2);
  \draw[force] ([shift={(110:0.6)}]0,0) arc (110:250:0.6);
  \node[left] at (-0.65,0) {$x_c^T F_y^T$};
\end{tikzpicture}\\[0.4em]
combined force + moment at $x=0$ (C.O.M.)
\end{minipage}

\bigskip
\noindent
$\rightarrow$ The Airbus A350 has a takeoff weight of \SI{600000}{lb} and symmetric wings.

\begin{center}
\begin{tikzpicture}[baseline=(b)]
  \coordinate (b) at (0,0);
  \draw[thick] (-3.4,0) -- (3.4,0);
  \com{(0,0)}
  \draw[force] (0,1.5) node[above] {$W$} -- (0,0.2);
  \draw[force] (2.6,-1.3) node[below] {$F_y^T$} -- (2.6,-0.03);
  \draw[force] (-2.6,-1.3) node[below] {$F_y^T$} -- (-2.6,-0.03);
  \draw[thinarrow] (0.25,0.75) -- (2.6,0.75) node[midway, above] {$x_c^T$};
  \draw[thinarrow] (-0.25,0.75) -- (-2.6,0.75) node[midway, above] {$x_c^T$};
\end{tikzpicture}
\quad{\Large $\equiv$}\quad
\begin{tikzpicture}[baseline=(b)]
  \coordinate (b) at (0,0);
  \com{(0,0)}
  \draw[force] (0,1.5) node[above] {$W$} -- (0,0.2);
  \draw[force] (0,-1.3) node[below] {$2F_y^T$} -- (0,-0.2);
  \draw[force] ([shift={(110:0.6)}]0,0) arc (110:250:0.6);
  \node[left] at (-0.65,0) {$x_c^T F_y^T$};
  \draw[force] ([shift={(70:0.6)}]0,0) arc (70:-70:0.6);
  \node[right] at (0.65,0) {$x_c^T F_y^T$};
\end{tikzpicture}
\end{center}

\begin{itemize}[label=$\rightarrow$]
  \item The moments cancel, so there is \emph{no plane rotation} (roll).
  \item $2F_y^T = \SI{630000}{lb} > W$, so $\sum F_y > 0$!
  \begin{itemize}[label=\then]
    \item The plane lifts off!
  \end{itemize}
\end{itemize}

\section*{Note for HW\,1: Statically Indeterminate Problems}

\begin{center}
\begin{tikzpicture}[baseline=(b), scale=0.8]
  \coordinate (b) at (0,0);
  \lwall{0}{-0.9}{0.9}
  \draw[beam] (0,-0.15) rectangle (5,0.15);
  \node[pin] at (0,0) {}; \node[left] at (-0.4,0) {$A$};
  \draw[force] (2.5,1.2) node[above] {$F$} -- (2.5,0.17);
  \pinsupport{5}{-0.15}
  \node[above] at (5,0.15) {$B$};
\end{tikzpicture}
\hspace{1cm}
\begin{tikzpicture}[baseline=(b), scale=0.8]
  \coordinate (b) at (0,0);
  \node[font=\bfseries] at (-0.8,1.6) {FBD};
  \draw[thick] (0,0) -- (5,0);
  \node[pin] at (0,0) {}; \node[pin] at (5,0) {};
  \draw[force] (-1.2,0) node[left] {$A_x$} -- (-0.08,0);
  \draw[force] (0,-1.1) node[below] {$A_y$} -- (0,-0.08);
  \draw[force] ([shift={(20:0.5)}]0,0) arc (20:160:0.5);
  \node at (0,0.85) {$M_A$};
  \draw[force] (2.5,1.2) node[above] {$F$} -- (2.5,0.08);
  \draw[force] (5.08,0) -- (6.1,0) node[right] {$B_x$};
  \draw[force] (5,-1.1) node[below] {$B_y$} -- (5,-0.08);
\end{tikzpicture}
\end{center}

\begin{itemize}
  \item Unknowns: $A_x,\ A_y,\ M_A,\ B_x,\ B_y$.
  \item Statics equations (2D planar): $\sum F_x = 0,\ \sum F_y = 0,\ \sum M = 0$.
\end{itemize}

\noindent
$\rightarrow$ If (\# unknowns) $>$ (\# independent equations):
\begin{itemize}[label=\then]
  \item We cannot solve the problem using statics alone!
  \item It is ``\emph{statically indeterminate}.''
  \item We need extra information on the \emph{beam deflection} to solve the problem.
        We will learn this later.
\end{itemize}

\end{document}
